\section{Introduction}
\label{sec:introduction}

\IEEEPARstart{T}{he} rapid digitization of clinical workflows, electronic health record (EHR) systems, and remote patient monitoring telemetry has highlighted the necessity for robust, computationally efficient data encryption schemes. Sensitive healthcare information—including patient identifiers, clinical diagnoses, prescribed medications, vital physiological measurements, and historical lab reports—requires strict confidentiality during transmission and storage. Furthermore, healthcare applications demand 100\% loss-less, exact reversible decryption so that decrypted diagnostic records remain completely uncorrupted.

Standard symmetric ciphers such as the Advanced Encryption Standard (AES) \cite{dworkin2007recommendation} provide strong cryptographic security. However, their internal designs often rely on complex substitution-permutation networks (SPN) and pre-computed S-boxes \cite{daemen2002design} that can impose memory and gate-count trade-offs in resource-constrained edge systems. Consequently, researchers have investigated alternative mathematical primitives capable of delivering fast non-linear bit manipulation and strong bit diffusion with lightweight software and hardware operations.

Cellular Automata (CA) present an attractive paradigm for lightweight cryptographic transformations due to their bit-level locality, parallelizable state updates, and simple logical operations \cite{wolfram1983statistical, wolfram1986random}. While early CA stream generators based on 1D Elementary Cellular Automata (such as Rule 30) suffered from cryptanalytic vulnerabilities under algebraic and structural attacks \cite{gutowitz1993cryptology, nandi2002theory}, multi-dimensional block-based CA structures offer improved structural diffusion when combined with spatial matrix shifts and reversible partition rules. In particular, Margolus neighborhood Cellular Automata provide an inherently reversible block transformation framework where local $2 \times 2$ matrix blocks evolve losslessly through bijective lookup rules.

To address the need for a practical, reversible block cryptosystem suited for healthcare payloads, this paper presents the \textbf{SHA-512-Based 256-bit Block Cellular Automata Cryptosystem}. The system derives a 256-bit binary key ($K_{256}$) directly from the first 32 bytes of a SHA-512 password digest. Input healthcare text is framed using a 32-bit big-endian binary length header, zero-padded to a 256-bit boundary, and split into 256-bit blocks mapped to $16 \times 16$ binary matrices. Non-linear state evolution is performed across multiple rounds of $2 \times 2$ Margolus CA transformations, alternating between standard partitions on odd rounds and shifted periodic partitions on even rounds, with a 2D circular matrix shift (down by 2 rows, right by 2 columns) applied after every round. Decryption reverses these steps in exact order, unwinding inverse circular shifts and inverse Margolus partitions before stripping padding to recover the exact UTF-8 text payload.

To determine an appropriate balance between non-linear diffusion and computational overhead, this study conducts a systematic round-scaling evaluation across:
\begin{equation}
R \in \{2, 4, 6, 8, 10\}
\end{equation}
The empirical investigation analyzes how the number of CA rounds influences changed-bit counts, bit avalanche percentage, Byte NPCR, Byte UACI, and execution latency. Furthermore, the complete pipeline is validated against 100 deterministic synthetic healthcare records to verify exact end-to-end record recovery. Integrity verification and authentication tags are outside the scope of this unauthenticated block cryptosystem and can be integrated via external MAC mechanisms in future work.

\subsection{Research Contributions}
\label{subsec:contributions}
The primary contributions of this work are summarized as follows:
\begin{enumerate}
    \item \textbf{SHA-512 Key Derivation Framework}: We establish a key derivation scheme that extracts a 256-bit binary key ($K_{256}$) from the first 32 bytes of a SHA-512 digest to drive 256-bit block XOR mixing.
    \item \textbf{Reversible 16x16 Binary Matrix Block Transformation}: We formalize a 256-bit block transformation pipeline utilizing $16 \times 16$ binary matrix representations, reversible $2 \times 2$ Margolus CA local rules, alternating standard/shifted periodic partitions, and 2D circular matrix shifts.
    \item \textbf{Reversible Payload Framing Scheme}: We implement a 32-bit binary length header framing mechanism with zero-padding to enable exact loss-less recovery of arbitrary UTF-8 healthcare text strings.
    \item \textbf{Systematic Round-Scaling Experiment}: We conduct a comprehensive empirical evaluation across $R \in \{2, 4, 6, 8, 10\}$ measuring bit avalanche percentage, Byte NPCR, Byte UACI, execution latency, reversibility, and healthcare record recovery.
    \item \textbf{Empirical 10-Round Configuration Selection}: Based on the measured diffusion/performance trade-off, we select 10 rounds as the proposed production configuration, achieving 81.92 changed bits (32.00\% avalanche), 84.62\% Byte NPCR, 22.83\% Byte UACI, 1.883 ms per-block latency, and 100\% healthcare record recovery.
\end{enumerate}

\subsection{Paper Organization}
\label{subsec:organization}
The remainder of this paper is organized as follows: Section \ref{sec:literature_review} reviews related literature on Cellular Automata and symmetric cryptography. Section \ref{sec:methodology} describes the mathematical formulation of the $2 \times 2$ Margolus CA rules and SHA-512 key derivation. Section \ref{sec:architecture} details the system architecture and end-to-end encryption/decryption workflows. Section \ref{sec:security_analysis} presents the security analysis and diffusion metrics. Section \ref{sec:benchmarks} evaluates experimental latency and performance scaling. Section \ref{sec:discussion} discusses structural trade-offs. Section \ref{sec:future_work} outlines future research directions, and Section \ref{sec:conclusion} concludes the paper.
